% L1b takeaway deck.
%
% THE DECK MIRRORS THE NOTEBOOK. The instructor shows
% CHEME-5660-L1b-Lecture-TimeValueMoney-Fall-2026.ipynb while students follow
% along and take notes on these slides. The deck therefore covers the same
% material in the same order, with frame titles matching notebook headings.
% Adding a section to the notebook means adding a slide here.
%
% Notebook outline as of this writing, and where each part lands:
%
%   (deck only)                          -> title
%   Disclaimer and Risks       [cell 8]  -> slide 2   (always near the front)
%   Learning Objectives        [cell 0]  -> Objectives for Today
%   Time Value of Money        [cell 1]  -> Time Value of Money
%   Discrete Compounding       [cell 2]  -> Discrete Compounding
%   Continuous Compounding     [cell 3]  -> Continuous Compounding
%   Who Sets the Valuation Rate[cell 4]  -> Who Sets the Valuation Rate? + What Rate Should I Use?
%   Abstract Assets            [cell 5]  -> An Abstract Asset (figure) + Net Cash Flows
%   Net Present Value (NPV)    [cell 6]  -> Net Present Value + Reading the Sign of NPV
%   Summary                    [cell 7]  -> Summary
%
% The lecture notebook carries no "Setup, Data, and Prerequisites" section and no
% code cells: L1b's code lives in the companion worked example and the optional
% derivation. There is nothing to skip, so the deck covers cells 0-8 in order,
% with cell 8 (the disclaimer) pulled to the front.
%
% Optional notebooks get NO frame of their own in this deck: the optional
% derivation, CHEME-5660-L1b-Derivation-Discrete-Continuous-Compounding-Fall-
% 2026.ipynb, is listed under "Optional:" on the Lectures and Examples frame and
% nowhere else. L3a/L3b/L4a carry an "Optional Advanced Material" frame instead;
% L1b deliberately does not.
%
% The 2024 decks opened with a "2-minute rapid review" transition frame. That is
% a holdover: it belongs to a concept-review notebook section, and L1b has none,
% so it is not used here. \transitionframe is still in vnslides.sty for decks
% whose notebook does open that way.
\documentclass[aspectratio=169,11pt]{beamer}
\usepackage{vnslides}

\title{{\vnheavy L1b}: Time Value of Money, Abstract Assets, and Net Present Value}
\subtitle{CHEME 5660 Cornell University}
\author{Copyright Jeffrey Varner 2026. All Rights Reserved}

\begin{document}

% ---- title ---------------------------------------------------
\vntitlepage

% ---- disclaimer ----------------------------------------------
\begin{frame}{Disclaimer and Risks}
  \vspace{0.6em}
  \textbf{\ink{Training only.}} This content is for training and informational
  purposes only. The teaching team does not make, give, or endorse any offer or
  solicitation to buy or sell securities or derivative products, or provide
  investment or trading advice or strategy.

  \vspace{1.1em}
  \textbf{\ink{Risk.}} Trading involves risk and is generally not recommended for
  those with limited resources, experience, or a low risk tolerance. Past
  performance does not guarantee future success.

  \vspace{1.1em}
  \textbf{\ink{Responsibility.}} You are responsible for your investment decisions
  based on your financial circumstances, objectives, risk tolerance, and
  liquidity needs.
\end{frame}

% ---- objectives ----------------------------------------------
\begin{frame}{Objectives for Today}
  \vspace{0.4em}
  Today we compare \hilite{cash flows at different dates}, derive accumulation
  and discount factors, and use them to value an \hilite{abstract asset}.

  \vspace{0.9em}
  \decklabel{Today's concepts}
  \vspace{0.4em}
  \begin{itemize}\setlength{\itemsep}{0.5em}
    \item \hilite{The time value of money}: cash flows at different dates cannot
      be combined until they are expressed at the same valuation date.
    \item \hilite{An abstract asset}: a sequence of dated signed cash flows. The
      same representation covers a bond, a stock, a factory, or a car.
    \item \hilite{Net present value}: discount each signed cash flow to a common
      valuation date, then add. The sign of the NPV is a measure of the attractiveness of an abstract asset relative to a declared discount rate.
  \end{itemize}
  \vspace{-0.4em}

\end{frame}

% ---- examples ------------------------------------------------
\begin{frame}{Lectures and Examples}
  \vspace{0.25em}
  {\normalsize\textbf{\ink{Lecture:}}\par
  \dimmed{CHEME-5660-L1b-Lecture-TimeValueMoney-Fall-2026.ipynb}}

  \vspace{0.8em}
  {\normalsize\textbf{\ink{Example:}}\par
  \dimmed{CHEME-5660-L1b-NetPresentValue-Fall-2026-WorkedExample.ipynb}}

  \vspace{0.5em}
  Represent ten years of Tesla Model S ownership as dated cash flows, discount
  them to time $0$, and compute their NPV.

  \vspace{0.8em}
  {\normalsize\textbf{\ink{Optional:}}\par
  \dimmed{CHEME-5660-L1b-Derivation-Discrete-Continuous-Compounding-Fall-2026.ipynb}}

  \vspace{0.5em}
  Measure how far discrete compounding falls short of the fixed-rate continuous
  limit, and when the estimate of that gap can be trusted.

  \slidenote{Link:}{\href{https://github.com/varnerlab/CHEME-5660-CourseRepository-Fall-2026.git}{https://github.com/varnerlab/CHEME-5660-CourseRepository-Fall-2026.git}}
\end{frame}

% ---- TVM -----------------------------------------------------
\begin{frame}{Time Value of Money}
  A dollar today and a dollar at a future date are \hilite{not comparable}. We must express both on a common
  valuation date using a conversion rule.

  \vspace{0.5em}
  Let $P$ be an amount held at time $0$, and let $y$ be a quoted
  \textbf{\ink{nominal annual yield}} used for valuation. With one (discrete) compounding
  period per year, the amount after one year is given by:
  \[
    F_1=\underbrace{(1+y)}_{\mathcal D_{1,0}(y;1)}P .
  \]

  For $y>0$ the factor $\mathcal D_{1,0}(y;1)>1$ (called the \textbf{\ink{accumulation factor}}) moves time-0 dollars forward in time.
  It's inverse, $\mathcal D_{1,0}^{-1}(y;1)<1$ (called the \textbf{\ink{discount factor}}) moves future dollars back in time.

  \slidenote{Careful:}{a later dollar is not \emph{automatically} worth less. That depends on the rate $y$ you declare.}
\end{frame}

% ---- math: accumulation --------------------------------------
\begin{frame}{Discrete Compounding}
  Let $j=0,1,2,\ldots$ index the \textbf{\ink{compounding periods}}, let
  $i_\ell>-1$ be the \mbox{\hilite{effective periodic interest rate}} over
  $\ell\rightarrow\ell+1$, and collect these as $\mathbf i=(i_0,\ldots,i_{j-1})$,
  empty at $j=0$. An amount $P$ at time $0$ becomes $F_j$ at period $j$:
  \vspace{-0.3em}
  \[
    F_j=\underbrace{\left[\prod_{\ell=0}^{j-1}(1+i_\ell)\right]}_{\mathcal D_{j,0}(\mathbf i)}P ,
    \qquad \text{\footnotesize empty product} = 1 .
  \]
  \vspace{-0.55em}

  For a nominal annual yield $y$ compounded $n$ times per year, every
  $i_\ell=y/n$:
  \vspace{-0.3em}
  \[
    \underbrace{\mathcal D_{j,0}(y;n)=\left(1+\frac{y}{n}\right)^{j}}_{\text{\footnotesize accumulation}},
    \qquad
    \underbrace{\mathcal D_{j,0}^{-1}(y;n)=\left(1+\frac{y}{n}\right)^{-j}}_{\text{\footnotesize discount}}.
  \]

  \vspace{0.05em}
  {\small \hilite{Assumes}: constant $y/n$, end-of-period compounding,
  reinvested interest, $1+y/n>0$.}

  \slidenote{Careful:}{a flat $y$ is an \emph{assumption}. In L2b, one security's yield gives way to a curve of spot rates.}
\end{frame}

% ---- math: continuous ----------------------------------------
\begin{frame}{Continuous Compounding}
  Let $g(t)$ be a \hilite{continuously compounded growth rate}, carrying units of
  inverse time, and let $T$ be the horizon in years. Over $0\rightarrow T$:
  \vspace{-0.2em}
  \[
    \mathcal D_{T,0}(g)=\exp\left(\int_0^T g(u)\,du\right),
    \qquad
    \mathcal D_{T,0}^{-1}(g)=\exp\left(-\int_0^T g(u)\,du\right),
  \]
  \vspace{-0.6em}

  which for \hilite{constant} $g$ collapses to:
  \vspace{-0.2em}
  \[
    \underbrace{\mathcal D_{T,0}(g)=e^{gT}}_{\text{\footnotesize accumulation}},
    \qquad
    \underbrace{\mathcal D_{T,0}^{-1}(g)=e^{-gT}}_{\text{\footnotesize discount}}.
  \]

  A nominal annual yield $y$ compounded $n$ times per year fits this form.
  Setting $g=g_y$ with $g_y:=n\log\left(1+y/n\right)$ makes
  $\left(1+y/n\right)^{nT}=e^{g_yT}$ hold exactly at every finite $n$: $y$ and
  its continuous equivalent $g_y$ are
  \textbf{\ink{two representations of one accumulation rule}}.

  \slidenote{Careful:}{at finite $n$, $y\neq g_y$, so $y$ does not belong in $e^{yT}$. Only as $n\rightarrow\infty$ does $g_y\rightarrow y$.}
\end{frame}

% ---- who sets the rate ---------------------------------------
\begin{frame}{Who Sets the Valuation Rate?}
  The valuation rate is \hilite{not universal}. It depends on the cash flow, market
  convention, horizon, and the risks being modeled: no single policy rate prices
  every asset.

  \vspace{0.5em}
  \eyebrow{The Federal Reserve}
  \vspace{0.25em}
  \begin{itemize}\setlength{\itemsep}{0.4em}
    \item The Fed conducts \textbf{\ink{monetary policy}} for maximum employment,
      stable prices, and moderate long-term rates, \textbf{\ink{supervises
      banks}}, and is the \textbf{\ink{lender of last resort}}.
    \item The Federal Open Market Committee (\textbf{\ink{FOMC}}) sets a target range for the federal funds rate, the
      overnight rate on unsecured reserve balances traded between eligible
      institutions. 
    \item Other market rates, e.g., Treasury yields and mortgage rates, respond to
      policy and market conditions, but are \textbf{\ink{not mechanically pegged}}
      to the federal funds rate.
  \end{itemize}

  \slidenote{Link:}{\href{https://www.federalreserve.gov/monetarypolicy/fomccalendars.htm}{FOMC calendar and meeting materials}}
\end{frame}

\begin{frame}{What Valuation Rate Should I Use?}
  There is no universal answer: $y$ is a \hilite{declared model input}, matched
  to the cash flow by currency, horizon, and any stated adjustment for credit,
  liquidity, or risk.

  \vspace{0.45em}
  \begin{center}\footnotesize
  \begin{tabular}{@{}>{\raggedright\arraybackslash}p{0.24\textwidth}c>{\raggedright\arraybackslash}p{0.53\textwidth}@{}}
    \toprule
    \textbf{\ink{Term}} & \textbf{\ink{Symbol}} & \textbf{\ink{Meaning here}}\\[1pt]
    \midrule
    \textbf{\ink{Interest rate}} & $i_\ell$ & A rate paid or received for
      borrowing or lending, such as the federal funds rate or a mortgage rate.
      Those are annual quotes; the model uses $i_\ell=y/n$ per period.\\
    \addlinespace[4pt]
    \textbf{\ink{Yield}} & $y$ & A rate tied to the price and cash flows of an
      investment, such as a bond yield.\\
    \addlinespace[4pt]
    \textbf{\ink{Discount rate}} or \textbf{\ink{valuation rate}} & $y$ & The rate
      a model uses to convert future cash flows to present value.\\
    \addlinespace[4pt]
    \textbf{\ink{Continuous equivalent}} & $g_y$ & The same rate written in
      continuous form, $g_y=n\log\left(1+y/n\right)$.\\[1pt]
    \bottomrule
  \end{tabular}
  \end{center}

  \slidenote{Careful:}{the same $y$ fills rows 2 and 3: one number, two roles. From L2a, $g$ is an asset's own measured growth rate.}
\end{frame}

% ---- figure --------------------------------------------------
\begin{frame}{An Abstract Asset}
  \vspace{0.2em}
  \begin{center}
    \resizebox{\textwidth}{!}{% Abstract-asset timeline -- SHARED SOURCE.
%
% \input by the standalone wrapper below; do not duplicate the picture.
% Rebuild the PDF/SVG for the notebook and slides with `make` in this directory.
%
% Drawing grammar and the CHEME flow convention come from templates/vnflow.sty:
%   arrow LEAVING a time-point unit  = cash paid
%   arrow ENTERING a time-point unit = cash received
% The signed net cash flow \bar c_j already carries its sign, so labels are
% written WITHOUT an extra minus; the arrow direction conveys direction.
%
% Indexing matches the notebooks and L2a: j = 0,1,...,N with N = nT.
\begin{tikzpicture}[x=2.25cm,y=1cm,>=Latex,font=\sffamily\small]

  \draw[vn brace] (0.35,1.65) -- (6.15,1.65)
    node[midway,above=7pt,text=vnink]
    {\bfseries abstract asset $\mathcal{A}=\{\bar c_0,\bar c_1,\ldots,\bar c_N\}$};

  \draw[vn timeline] (0,0) -- (6.65,0) node[right] {time};
  \foreach \x/\lab/\unit in {0.35/0/u0,1.75/1/u1,3.15/2/u2,6.0/N/uN}{
    \node[vn unit] (\unit) at (\x,0) {};
    \node[below=4pt] at (\x,0) {$j=\lab$};
  }
  \node at (4.55,0.02) {$\cdots$};

  % j = 0: arrow leaves the unit -> cash paid (the purchase)
  \draw[vn paid] (u0.north) -- (0.35,1.05);
  \node[right,text=vncarnelian] at (0.35,0.62) {$\bar c_0$};
  \node[above,vn annotation] at (0.35,1.05) {cash paid};

  % j = 1: arrow enters the unit -> cash received
  \draw[vn received] (1.75,1.05) -- (u1.north);
  \node[right,text=vnlink] at (1.75,0.62) {$\bar c_1$};

  % j = 2: arrow leaves the unit -> cash paid
  \draw[vn paid] (u2.north) -- (3.15,1.05);
  \node[right,text=vncarnelian] at (3.15,0.62) {$\bar c_2$};

  % j = N: arrow enters the unit -> cash received (the sale)
  \draw[vn received] (6.0,1.05) -- (uN.north);
  \node[right,text=vnlink] at (6.0,0.62) {$\bar c_N$};
  \node[above,vn annotation] at (6.0,1.05) {cash received};

  % Transport every dated flow to the common valuation date
  \draw[vn transport] (6.0,-0.94) -- (0.50,-0.94)
    node[midway,fill=white,inner sep=3pt,text=vnlink]
    {$\bar c_j\mapsto\mathcal{D}_{j,0}^{-1}(y;n)\,\bar c_j$};
  \node[vn callout,anchor=north west] at (0.35,-1.27)
    {common units: time-0 dollars};
  \node[vn annotation,anchor=west] at (3.35,-1.57)
    {discount future cash flows to the valuation date};
\end{tikzpicture}
}
  \end{center}
\end{frame}

% ---- net cash flows ------------------------------------------
\begin{frame}{Net Cash Flows}
  An abstract asset is a \textbf{\ink{sequence of dated cash flows}} in one
  currency, one at each compounding date $j=0,1,\ldots,N$, so $N$ is the final
  index and there are $N+1$ dates. At date $j$, let
  $\mathbf c_j\in\mathbb R_{\geq0}^{m_j}$ hold the $m_j$
  nonnegative cash flow magnitudes and $\symbf{\nu}_j\in\{-1,+1\}^{m_j}$ their directions,
  $+1$ an inflow and $-1$ an outflow. The \hilite{signed net cash flow} is their
  inner product:
  \[
    \bar c_j=\left\langle\mathbf c_j,\symbf{\nu}_j\right\rangle
    =\sum_{k=1}^{m_j}c_{j,k}\,\nu_{j,k}.
  \]

  The undiscounted total $\sum_{j=0}^{N}\bar c_j$ is \textbf{\ink{not}} a time-0
  value: its terms are dollars from different dates.

  \slidenote{Next:}{discount every term to a common date, and the result is net present value.}
\end{frame}

% ---- math: NPV -----------------------------------------------
\begin{frame}{Net Present Value}
  \vspace{0.2em}
  Take a nominal annual yield $y$ compounded $n$ times per year as the declared
  valuation rate. Discount every signed net cash flow $\bar c_j$ to the valuation
  date with $\mathcal D_{j,0}^{-1}(y;n)$, \hilite{then add}. The net present value
  at time $0$ is given by:
  \[
    \operatorname{NPV}_0(y)=\sum_{j=0}^{N}\mathcal D_{j,0}^{-1}(y;n)\,\bar c_j
    =\underbrace{\left\langle \mathcal D_{\star,0}^{-1}(y;n),\ \bar{\mathbf c}\right\rangle}_{\text{\footnotesize time-0 dollars}} .
  \]

  \vspace{0.2em}
  The $\star$ runs over every cash-flow index, so $\mathcal D_{\star,0}^{-1}(y;n)$
  and $\bar{\mathbf c}$ are ordered $(N+1)$-vectors and NPV is a single scalar
  product, which is exactly how we compute it in code.

  \slidenote{Worked example:}{ten years of Model S ownership, $\operatorname{NPV}_0\approx-\$121{,}484$}
\end{frame}

% ---- NPV meaning ---------------------------------------------
\begin{frame}{Reading the Sign of NPV}
  \vspace{0.2em}
  The sign of $\operatorname{NPV}_0(y)$ compares the modeled cash flows against
  the \hilite{rate you declared}. It is a valuation statement, not a promise of
  accounting profit or realized return.

  \vspace{0.6em}
  \begin{itemize}\setlength{\itemsep}{0.45em}
    \item \textbf{\ink{Positive}}: in time-0 dollars, receipts exceed payments,
      so at the declared $y$ the asset is worth more than its price.
    \item \textbf{\ink{Negative}}: receipts fall short of payments, so at the
      declared $y$ the asset is worth less than its price.
    \item \textbf{\ink{Zero}}: receipts and payments agree, so the asset is worth
      exactly its price. This is how L2a prices Treasury bills, notes, and bonds.
  \end{itemize}

  \vspace{0.6em}
  Change $y$ and the sign can change: it describes this cash-flow model at this
  rate, not the asset in general.

  \slidenote{Careful:}{a positive NPV is not a claim that an asset is risk-free or universally attractive.}
\end{frame}

% ---- summary -------------------------------------------------
% Summary frames follow the notebook's Summary cell: one summary sentence, the
% three key takeaways, one concluding sentence. See vnslides.sty rule 7.
\begin{frame}{Summary}
  This lecture introduced time-value-of-money calculations and the
  representation of investments as dated cash flows.

  \vspace{0.6em}
  \decklabel{Key takeaways}
  \vspace{0.25em}
  \begin{itemize}\setlength{\itemsep}{0.45em}
    \item \textbf{\ink{Dated dollars must be compared at a common valuation date.}}
      A forward accumulation factor carries value forward and its inverse
      discounts it back, so cash flows can be added only once they share a date.
    \item \textbf{\ink{Rate conventions must be stated.}} A nominal yield quoted
      with a compounding frequency and its continuously compounded equivalent
      describe one accumulation rule, so the convention travels with the number.
    \item \textbf{\ink{NPV values an asset relative to a declared rate.}} The sign
      of the result is a statement about the cash-flow model and the discount
      rate you declared.
  \end{itemize}

  \vspace{0.6em}
  These ideas provide the valuation framework used throughout the course.
\end{frame}

\end{document}
